\begin{table}[htbp]\centering\footnotesize
\caption{Gasto federalizado: participaciones y los ocho fondos del Ramo 33, línea P (mdp)}\label{cua:feder}
\begin{tabular}{lrrr}
\toprule
Concepto & 2026 & 2027 & Var. real (\%) \\
\midrule
Ramo 28 Participaciones (no condicionadas) & 1.456.045,9 & 1.547.607,9 & +2,2 \\
\midrule \multicolumn{4}{l}{\emph{Ramo 33, aportaciones (condicionadas)}} \\
\quad FONE, nómina educativa & 546.396,8 & 606.878,2 & +6,8 \\
\quad FASSA, servicios de salud & 84.635,9 & 89.714,0 & +1,9 \\
\quad FAIS, infraestructura social & 135.060,7 & 144.751,3 & +3,0 \\
\quad FORTAMUN, municipios & 136.817,4 & 146.634,1 & +3,0 \\
\quad FAM, aportaciones múltiples & 43.464,6 & 46.583,2 & +3,0 \\
\quad FAETA, educación tecnológica & 10.811,4 & 11.224,3 & $-$0,2 \\
\quad FASP, seguridad pública & 9.951,1 & 10.447,9 & +0,9 \\
\quad FAFEF, entidades federativas & 74.754,9 & 80.118,5 & +3,0 \\
\textbf{Suma de los ocho fondos} & \textbf{1.041.892,8} & \textbf{1.136.351,5} & \textbf{+4,9} \\
Ramo 33 según el Anexo 1 & 1.041.892,9 & 1.136.351,6 & +4,9 \\
\bottomrule
\end{tabular}
\\[2pt]\begin{minipage}{0.92\textwidth}\scriptsize \textbf{Los ocho fondos suman el total del Ramo 33 al peso en los dos ejercicios}, lo que valida la lectura del Anexo 22 contra el Anexo 1. Participaciones y aportaciones \textbf{no se suman entre sí}: las primeras no están condicionadas y las segundas sí.\end{minipage}
\end{table}
